\documentclass[11pt]{article}
\usepackage[T1]{fontenc}
\usepackage{lmodern}
\usepackage{amsmath,amssymb,amsthm,mathtools}
\usepackage{microtype}
\usepackage{geometry}
\usepackage{hyperref}
\usepackage{enumitem}
\geometry{margin=1in}

\newtheorem{theorem}{Theorem}[section]
\newtheorem{lemma}[theorem]{Lemma}
\newtheorem{proposition}[theorem]{Proposition}
\theoremstyle{definition}
\newtheorem{definition}[theorem]{Definition}
\theoremstyle{remark}
\newtheorem{remark}[theorem]{Remark}

\title{Degree-Four Pachner Sites and Target-Present Five-Tetrahedron Closure}
\author{Inacio F. Vasquez\\Independent Researcher, United States}
\date{}

\begin{document}
\maketitle

\begin{abstract}
We isolate a local structural result in a constructive program based on
Pachner moves on closed triangulations.  For a degree-four vertex with
connected represented vertex link, the formal development constructs a
$4\to1$ site and proves a dichotomy: the site is legal, or its opposite
target tetrahedron is represented.  In the latter branch, the represented
target together with the four source tetrahedra forms a pairwise distinct
five-tetrahedron configuration, and the configuration closes the
represented vertex link.  The statements are formalized in Lean 4.

The result is deliberately local.  It does not prove a global strict
descent theorem and does not claim a proof of the Poincare conjecture.
\end{abstract}

\section{Introduction}

A constructive approach to triangulated $3$-manifolds is to analyze local
Pachner moves and seek a complexity-decreasing branch.  A recurring
difficulty is the $4\to1$ configuration at a degree-four vertex: the
candidate site may fail to be immediately legal because its opposite
tetrahedron is already represented.

This paper isolates the formal result that resolves that local branch into
a precise alternative.  The result is useful independently of the
remaining global descent problem.

The formal development is hosted at
\url{https://github.com/inaciovasquez2020/poincare-new-derivation}.

\section{Combinatorial setting}

Let $K$ be a finite triangulation whose represented tetrahedra are
\[
  K.\mathrm{tets}.
\]
A tetrahedron is written $\tau$, with vertex list $\tau.\mathrm{verts}$.
For a vertex $v$, let $\operatorname{vertexDegree}(K,v)$ denote the
represented tetrahedron degree used by the formal development.

A \emph{closed triangulation core} is the predicate
\texttt{ClosedTriangulationCore K}.  Its incidence component states that a
represented nondegenerate triangular face occurs in exactly two represented
tetrahedra.

For a vertex $v$, the represented vertex-link triangles are
\[
  \operatorname{vertexLinkTriangles}(K,v),
\]
and \texttt{VertexLinkConnected K v} is the formal connectivity condition
used below.

\section{The $4\to1$ site}

A \texttt{Move41Site} $s$ has center $s.e$, four source tetrahedra
\[
  s.\mathrm{sourceTet}_0,\ldots,s.\mathrm{sourceTet}_3,
\]
and an opposite target tetrahedron $s.\mathrm{targetTet}$.

Up to the site's vertex labels, the five model tetrahedra have the form
\[
\begin{aligned}
S_0&=\{a,b,c,e\},&
S_1&=\{a,b,d,e\},\\
S_2&=\{a,c,d,e\},&
S_3&=\{b,c,d,e\},\\
T&=\{a,b,c,d\}.
\end{aligned}
\]

The formal predicate \texttt{SameTetVertices $\tau$ $\sigma$} identifies
tetrahedra with the same represented vertex set.

\section{Incidence-two forcing}

\begin{lemma}[Incidence-two forcing]
Let $K$ satisfy \texttt{ClosedTriangulationCore K}.  Let $\tau\ne\rho$
be represented tetrahedra containing the same represented nondegenerate
triangular face.  If $\sigma$ is any represented tetrahedron containing
that face, then
\[
  \sigma=\tau\quad\text{or}\quad\sigma=\rho.
\]
\end{lemma}

\begin{proof}
The closed-core incidence condition says that the filtered list of
represented tetrahedra containing the face has length exactly two.
Since $\tau$ and $\rho$ are distinct members of that list, they exhaust
its possible members.
\end{proof}

The Lean declaration is
\begin{center}
\texttt{ClosedTriangulationCore.eq\_left\_or\_eq\_right\_of\_common\_face}.
\end{center}

\section{The target-present cluster}

Assume that all four source vertex sets occur exactly once in $K$ and that
the target vertex set is represented.  The formal development constructs
represented tetrahedra
\[
  \tau_0,\tau_1,\tau_2,\tau_3,\tau_T
\]
matching the four sources and the target.

It further proves
\[
  [\tau_0,\tau_1,\tau_2,\tau_3,\tau_T].\mathrm{Nodup}.
\]

The corresponding Lean declarations are
\begin{center}
\texttt{Move41Site.exists\_represented\_fiveTetCluster\_of\_targetPresent},
\end{center}
and
\begin{center}
\texttt{Move41Site.exists\_represented\_fiveTetCluster\_nodup\_of\_targetPresent}.
\end{center}

\begin{proposition}[Finite common-face closure]
For the represented target-present cluster above, every represented
tetrahedron containing a nondegenerate triangular face of a cluster member
is itself one of the five cluster members.
\end{proposition}

\begin{proof}
Each triangular face of each model tetrahedron has its required partner
among the other four model tetrahedra.  The incidence-two lemma then
excludes a third represented tetrahedron across that face.
\end{proof}

The exact formal theorem is
\begin{center}
\texttt{Move41Site.represented\_fiveTetCluster\_closed\_under\_common\_face}.
\end{center}

\section{Propagation through connected vertex links}

The local face-closure result propagates through a connected represented
vertex link.

\begin{lemma}[Link propagation]
Let $v$ have a connected represented vertex link.  If a nonempty family of
represented link triangles is closed under link adjacency, then it
contains every represented link triangle.
\end{lemma}

\begin{proof}
Connectivity gives a finite adjacency path from a chosen seed triangle to
any represented link triangle.  Induction along that path transports the
closure property.
\end{proof}

The formal declaration is
\begin{center}
\texttt{VertexLinkConnected.all\_of\_adjacent\_closed}.
\end{center}

Combining this propagation with common-face closure gives the formal
vertex-star saturation theorem
\begin{center}
\texttt{ClosedTriangulationCore.tetrahedronCluster\_vertexStar\_closed\_of\_commonFace\_closed}.
\end{center}

For the concrete target-present cluster, the resulting theorem is
\begin{center}
\texttt{Move41Site.represented\_fiveTetCluster\_vertexStarClosed\_of\_connectedLinks}.
\end{center}

It assumes a closed core, connected supported vertex links, the five
represented tetrahedra, their vertex-set identifications, and the
five-element nodup certificate, and concludes that every represented
tetrahedron in the vertex star of a cluster member is again one of the
five cluster members.

\section{Main theorem}

The principal result of this paper is the following exact formal
statement.

\begin{theorem}[Degree-four legal-or-target-present dichotomy]
Let $K$ satisfy \texttt{ClosedTriangulationCore K}.  Let $v$ satisfy
\[
  \operatorname{vertexDegree}(K,v)=4
\]
and assume
\texttt{VertexLinkConnected K v}.  Then there exists a
\texttt{Move41Site} $s$ such that
\[
  s.e=v
\]
and either
\[
  s.\mathrm{LegalIn}(K),
\]
or both of the following hold:
\begin{enumerate}[label=(\roman*)]
  \item there exists a represented tetrahedron $\tau$ with
  \texttt{SameTetVertices $\tau$ $s.targetTet$};
  \item every represented vertex-link triangle at $v$ is represented by
  a tetrahedron whose vertex set is one of the four source tetrahedra of
  $s$.
\end{enumerate}
\end{theorem}

The exact Lean declaration is
\begin{center}
\texttt{ClosedTriangulationCore.exists\_move41Site\_legalIn\_or\_targetPresent\_closes\_vertexLink}.
\end{center}

Its conclusion is, verbatim at the logical level,
\[
\exists s:\mathrm{Move41Site},\quad
s.e=v\ \land\left(
s.\mathrm{LegalIn}(K)
\ \lor\left[
\begin{array}{l}
\exists\tau\in K.\mathrm{tets},\mathrm{SameTetVertices}(\tau,s.\mathrm{targetTet})\\
\land\
\forall\sigma\in\operatorname{vertexLinkTriangles}(K,v),\exists\tau\in K.\mathrm{tets},\tau.\mathrm{linkTriangleAt?}\ v=\mathrm{some}\ \sigma\\
\qquad\land\exists\mathrm{source}\in s.\mathrm{sourceTets},\mathrm{SameTetVertices}(\tau,\mathrm{source})
\end{array}
\right]
\right).
\]

\begin{proof}
The degree-four hypothesis constructs the Move41 site and labels its
outer vertices.  The same degree hypothesis gives exact occurrence of
each of the four source vertex sets.

The outer-vertex coverage then shows that every represented tetrahedron
incident to the center is represented by one of the four source vertex
sets.  The source configurations are pairwise distinct.

The formal Move41 classification gives two alternatives.  In the first,
the site is legal.  In the second, the opposite target is represented.
The target-present branch is then combined with the five-tetrahedron
closure and connected-link propagation theorem, yielding the stated
vertex-link coverage.
\end{proof}

\section{Relation to global descent}

The theorem is a local structural result.  It does not imply that every
nonterminal triangulation admits a globally strict decrease of a
complexity functional.

In particular, this paper does not prove
\[
  K\not\simeq S^3
  \Longrightarrow
  \exists K'\;[\Phi(K')<\Phi(K)]
\]
for an arbitrary triangulation $K$, nor does it prove
\[
  \Phi(K)=0\Longrightarrow K\simeq S^3.
\]

Consequently, no proof of the Poincare conjecture is claimed here.

\section{Formal verification record}

The repository is
\url{https://github.com/inaciovasquez2020/poincare-new-derivation}.

The manuscript is synchronized to repository main commit
\texttt{dc517e7146538c69548301371e1d17fa1b8221e8}.

The repository's Lean toolchain is pinned by \texttt{lean-toolchain}; the
current value is
\texttt{leanprover/lean4:v4.30.0-rc1}.

The principal source modules are:
\begin{itemize}
  \item \texttt{lean/Poincare/DegreeFourConnectedLinkClassification.lean};
  \item \texttt{lean/Poincare/GlobalDegreeFourTargetPresentSaturation.lean}.
\end{itemize}

The manuscript should be considered a mathematical description of those
formal declarations, not a claim that the repository's entire research
program is closed.

\section{Conclusion}

The degree-four branch has a precise local structure.  A degree-four
vertex with connected represented link admits a Move41 site which is either
legal or enters a represented-target branch.  The latter branch produces a
five-tetrahedron configuration with exact incidence closure and vertex-link
saturation.

This isolates a concrete, machine-checked local theorem from the larger
global descent program while preserving the unresolved global boundary.

\end{document}
